% !PRINTED 640
\weberpage{640}
%<<CONT
und folglich kann $\frac{1}{2}(\alpha \pm \beta)$ keine ganze Zahl sein, weil der
absolute Werth der Norm dieser Zahl kleiner als 1 ist. Wenn
nun $\omega_1, \omega_2, \dots, \omega_n$ eine Basis von $\mathfrak{o}$ ist, und
\[
\begin{aligned}
\alpha &= a_1\omega_1 + a_2\omega_2 + \dots + a_n\omega_n \\
\beta &= b_1\omega_1 + b_2\omega_2 + \dots + b_n\omega_n,
\end{aligned}
\]
so können die ganzen rationalen Zahlen $a, b$ nicht den Congruenzen
\[ a_1 \equiv b_1,\ a_2 \equiv b_2,\ \dots,\ a_n \equiv b_n \pmod{2} \]
genügen, weil sonst $\frac{1}{2}(\alpha - \beta)$ eine ganze Zahl wäre.

Wenn wir also die sämmtlichen Zahlen $\alpha$ in $2^n$ Fächer vertheilen,
indem wir alle Zahlen in ein Fach werfen, in denen
$a_1, a_2, \dots, a_n$ dieselben Reste (0 oder 1) nach dem Modul 2
lassen, so können in jedem dieser Fächer höchstens 2 Zahlen,
nämlich $\alpha$ und $-\alpha$, vorkommen, und wenn wir $\alpha = 0$ noch
ausschliessen, so giebt es sogar in einem dieser Fächer, in dem
die $a_1, a_2, \dots, a_n$ alle gerade sind, gar keine Zahl $\alpha$. Die Anzahl
 aller möglichen Zahlen $\alpha$ ist also endlich und höchstens
$= 2^{n+1} - 2$.

Wenn nun der absolute Werth von $\alpha = 1$ ist, so gilt dasselbe
von allen Potenzen von $\alpha$. Folglich muss in der unbegrenzten
Reihe
\[ 1,\ \alpha,\ \alpha^2,\ \alpha^3,\ \dots \]
nothwendig dieselbe Zahl zum zweiten Male wiederkehren, also
$\alpha^h = \alpha^k$ und $k > h$ sein. Dann ist aber
\[ \alpha^{k-h} = 1, \]
d.\,h. $\alpha$ ist eine Einheitswurzel, wie bewiesen werden sollte.

Wir fügen noch die Bemerkung bei, dass, wenn eine Zahl $\alpha$
eines Körpers $\Omega$ eine Einheitswurzel $m^{\text{ten}}$ Grades ist, auch alle
mit $\alpha$ conjugirten Zahlen Einheitswurzeln desselben Grades sind.
Denn in der rationalen Gleichung $\alpha^m - 1 = 0$ kann $\alpha$ durch
jeden der conjugirten Werthe ersetzt werden.

Die Anzahl der Einheitswurzeln, die in einem Körper $\Omega$
enthalten sind, kann immer nur eine endliche sein, weil jede
Zahl in $\Omega$ einer rationalen Gleichung genügt, deren Grad dem
Körpergrad höchstens gleich ist. Der Grad der Einheitswurzeln
in $\Omega$ kann daher einen endlichen Werth nicht übersteigen.
